\chapter*{Anexo D: Resolución de identidad en detalle}
\addcontentsline{toc}{chapter}{Anexo D: Resolución de identidad en detalle}
\label{cap:anexo_resolucion}

\setcounter{section}{0}
\renewcommand{\thesection}{D.\arabic{section}}

Este anexo amplía la Sección~\ref{sec:procesamiento_resolucion} con el
detalle de las reglas, el análisis de errores que llevó a fijar su diseño
final y la sensibilidad de los resultados a sus parámetros. Todas las cifras
proceden de la ejecución de 10\,000 documentos CVN descrita en el
Capítulo~\ref{cap:procesamiento}, evaluada contra el fichero de control que
genera el propio generador sintético.

\section{Normalización previa a la comparación}
\label{sec:anexo_resolucion_normalizacion}

La comparación no opera sobre los textos originales sino sobre formas
normalizadas, que se guardan en silver junto a ellos.

\begin{itemize}
  \item \textbf{Nombres.} Se pliegan sin acentos, en minúsculas y sin
  puntuación. Las partículas de apellido, como «de», «del» o «la», se omiten
  al elegir la clave de bloqueo, que combina el primer apellido significativo
  y la inicial del nombre.
  \item \textbf{Organizaciones.} Se reducen a su conjunto ordenado de palabras
  significativas, sin las de relleno y sin las direcciones web que algunos
  perfiles incluyen en el nombre. El orden de las palabras deja de importar
  y se puede calcular un solape de conjuntos.
  \item \textbf{Obras.} El DOI pierde sus prefijos y pasa a minúsculas, y el
  título se normaliza igual que los nombres.
\end{itemize}

\section{Compatibilidad de nombres}
\label{sec:anexo_resolucion_nombres}

La regla de repliegue acepta una coincidencia de nombre en tres formas, que
se clasifican y se guardan en la evidencia de cada enlace.

\begin{itemize}
  \item \textbf{Completa.} Apellido y nombre coinciden tras la normalización.
  \item \textbf{Nombre relajado.} El apellido coincide y el nombre difiere en
  que una parte es solo una inicial o en que un nombre compuesto aparece
  abreviado, como «Ana M» frente a «Ana María». Un nombre que pierde una
  parte entera sigue sin ser compatible.
  \item \textbf{Apellido acortado.} El nombre coincide y un apellido es el
  primer apellido del otro, porque muchas personas omiten el segundo.
\end{itemize}

Un nombre que difiere por ambos lados a la vez no es compatible. Esta
restricción nació del análisis de errores de la
Sección~\ref{sec:anexo_resolucion_errores}.

\section{Evidencia guardada en cada enlace}
\label{sec:anexo_resolucion_evidencia}

Cada fila de la tabla de enlaces conserva la regla aplicada y una evidencia en
JSON con la forma de coincidencia de nombre y las organizaciones compartidas.
Dos indicadores marcan los casos dudosos. El de ambigüedad se activa cuando
un registro sin identificador tiene más de una entidad candidata, y en ese
caso queda como entidad propia. El de conflicto de nombre se activa cuando
un CVN cuyo nombre no es compatible con el de ninguno de los registros de
ORCID que comparten su identificador se fusiona con ellos. La fusión se
mantiene, porque el identificador es la clave autoritativa, y el indicador
permite revisarla.

\section{Umbral de similitud entre organizaciones}
\label{sec:anexo_resolucion_umbral}

La regla de repliegue exige que alguna organización coincida. El umbral por
defecto exige igualdad de las formas normalizadas, y se fijó con datos
reales y no por preferencia. Un solape de conjuntos de palabras no separa
bien los casos. La misma universidad con un sufijo de campus o de facultad
puntúa entre 0,40 y 0,60, mientras que dos universidades distintas de la
misma ciudad puntúan 0,67. Ningún umbral intermedio acepta lo primero y
rechaza lo segundo, y las traducciones de una misma institución entre
español e inglés puntúan 0,00 y no se detectan. La
Tabla~\ref{tab:anexo_umbral} muestra el efecto de bajarlo sobre la ejecución
de referencia.

\begin{table}[H]
  \centering
  \small
  \renewcommand{\arraystretch}{1.15}
  \begin{tabular}{p{0.22\textwidth} p{0.14\textwidth} p{0.14\textwidth} p{0.14\textwidth} p{0.14\textwidth}}
    \hline
    \textbf{Umbral} & \textbf{Fusiones} & \textbf{Correctas} & \textbf{Falsas} & \textbf{Ambiguas} \\
    \hline
    1,0 a 0,6 & 144 & 138 & 6 & 0 \\
    0,5 & 147 & 136 & 11 & 2 \\
    0,34 & 149 & 136 & 13 & 2 \\
    \hline
  \end{tabular}
  \caption{Efecto del umbral de similitud entre organizaciones.}
  \label{tab:anexo_umbral}
\end{table}

\section{Análisis de errores y decisión final}
\label{sec:anexo_resolucion_errores}

La primera medición a gran escala dio una precisión de la regla de repliegue
del 90,8\%, por debajo del objetivo del 99\% fijado antes de medir. Antes de
cambiar nada se analizaron los casos, y el análisis encontró dos efectos
distintos.

\begin{itemize}
  \item \textbf{Una evaluación demasiado estricta.} Se consideraba evaluable
  un documento solo si su semilla tenía un registro de ORCID, de modo que una
  fusión correcta con una entidad anclada por otro CVN que declaraba el mismo
  identificador contaba como fusión sin contraparte. Se corrigió la
  definición y los documentos evaluables pasaron de 180 a 186.
  \item \textbf{Una combinación de evidencia débil.} Los 14 falsos positivos
  compartían exactamente una organización, mientras que las 78 fusiones con dos
  organizaciones compartidas o más fueron correctas sin excepción. Por tipo
  de coincidencia, la de nombre completo dio 79 aciertos y 2 errores, la de
  nombre relajado 47 y 0, la de apellido acortado 12 y 4, y la que relajaba
  ambas partes a la vez 0 aciertos y 8 errores.
\end{itemize}

Prohibir relajar ambas partes del nombre a la vez eliminó los 8 errores sin
perder ningún acierto, y es la decisión que quedó en el diseño final. El
resultado fue una precisión del 95,8\% y una cobertura del 74,2\%, medida en
los 186 documentos evaluables.

\section{Alternativas más estrictas}
\label{sec:anexo_resolucion_alternativas}

Se midieron variantes que exigen más evidencia. Ninguna alcanza el 99\% y
todas cuestan cobertura, por lo que no se adoptaron por defecto. La
evidencia de cada enlace permite aplicarlas en una consulta posterior sin
repetir la ejecución. La Tabla~\ref{tab:anexo_estrictas} las compara.

\begin{table}[H]
  \centering
  \small
  \renewcommand{\arraystretch}{1.15}
  \begin{tabular}{
    >{\raggedright\arraybackslash}p{0.44\textwidth}
    p{0.14\textwidth} p{0.14\textwidth}}
    \hline
    \textbf{Exigencia adicional} & \textbf{Precisión} & \textbf{Cobertura} \\
    \hline
    Ninguna, diseño adoptado & 95,8\% & 74,2\% \\
    Dos organizaciones si el apellido está acortado & 98,5\% & 70,4\% \\
    Dos organizaciones si el nombre no es idéntico & 98,3\% & 62,9\% \\
    Dos organizaciones en cualquier fusión & 100\% & 41,9\% \\
    \hline
  \end{tabular}
  \caption{Variantes más estrictas de la regla de repliegue.}
  \label{tab:anexo_estrictas}
\end{table}

\section{Cobertura por variante de nombre}
\label{sec:anexo_resolucion_variantes}

El generador sintético escribe el nombre de los documentos sin identificador
como nombre exacto o con una de cuatro variantes, y la Tabla~\ref{tab:anexo_variantes} da la cobertura de la
regla en cada una. Los documentos que quedan sin fusionar son sobre todo los
que no declaran ninguna afiliación, 47 de los 186 evaluables, para los que la
regla no tiene evidencia por diseño. Esa condición fija un techo del 74,7\%
a la cobertura alcanzable, y la regla alcanza el 99,3\% de ese techo.

\begin{table}[H]
  \centering
  \small
  \renewcommand{\arraystretch}{1.15}
  \begin{tabular}{
    >{\raggedright\arraybackslash}p{0.4\textwidth}
    p{0.2\textwidth} p{0.2\textwidth}}
    \hline
    \textbf{Variante de nombre} & \textbf{Fusionados} & \textbf{Evaluables} \\
    \hline
    Nombre exacto & 43 & 59 \\
    Inicial del nombre & 47 & 60 \\
    Apellido en mayúsculas & 26 & 31 \\
    Solo el primer apellido & 12 & 19 \\
    Sin acentos & 10 & 17 \\
    \hline
  \end{tabular}
  \caption{Cobertura de la regla de repliegue por variante de nombre.}
  \label{tab:anexo_variantes}
\end{table}

\section{Método de evaluación}
\label{sec:anexo_resolucion_evaluacion}

El generador escribe, para cada documento, el identificador ORCID de su
semilla y el tipo de enlace que se le dio. Con ese fichero se conoce la
entidad correcta de cada documento. La evaluación cuenta como evaluable todo
documento cuya semilla tiene algún registro en silver, ya sea de ORCID o de otro CVN
que declara el mismo identificador. Una fusión de un
documento cuya semilla no está presente cuenta como falsa, porque nada
correcto podía fusionarse con él, y es la medida directa de que los
registros sin relación permanecen separados. Con la ejecución por defecto de 1\,000 documentos solo 16 de 288 tenían contraparte, frente a los 186 de la ejecución de 10\,000, porque el solapamiento entre las dos fuentes depende de cuánto de ORCID se aterrizó. Este método mide la
resolución contra los datos sintéticos de la plataforma y no contra datos
reales anotados, una limitación recogida en el
Anexo~\ref{cap:anexo_limitaciones}.
